\section{Decisions to confirm}
\label{sec:decisions}

Each box gives the recommendation, the reason and the alternative that was weighed. The team
delegated Assignment 2's open decisions to Cao Tri Tran on 23 October; no such delegation is
recorded for Assignment 3, so these are proposals until the Week 9 meeting records them. D-A3-2 to
D-A3-4 are rechecked when the Week 9--12 slides are released.

\begin{decision}{D-A3-1: a new, private repository with Assignment 2's runtime copied in}
\textbf{Recommendation.} \code{COS30049A3} holds \fp{backend/} and \fp{frontend/}. A script copies
the modules, models and snapshot tables that prediction needs out of the Assignment 2 repository at
the submitted commit, and records that commit and a SHA-256 for every file
(Section~\ref{sec:vendoring}). The repository stays private until marks are released.

\textbf{Why.} The brief wants the model code inside the back end's zip. The Assignment 2 repository
is over 200 MB with training code, and it is under marking, so it should not change. The saved
models are pickles that name their classes by module path, so the package must still be called
\code{explorer}; copying whole modules unchanged keeps them loadable.

\textbf{Alternative.} Adding \fp{backend/} and \fp{frontend/} to the Assignment 2 repository:
simpler imports, but it edits a submitted repository and ships data no web feature uses.
\end{decision}

\begin{decision}{D-A3-2: the front end grows from the Assignment 1 prototype}
\textbf{Recommendation.} Keep React, TypeScript, Vite, the prototype's own design system, its hash
router and its tests, and add to them.

\textbf{Why.} Assignment 1 promised that the screens would move to real data without being
rewritten, and the prototype already has the states, the validation, the responsive layouts and the
accessibility work that two usability rows of the rubric ask for.

\textbf{Alternative.} Starting again with create-react-app and JavaScript as in Workshop 8.
TypeScript is a superset of JavaScript and the README's commands are the same (\code{npm install},
\code{npm run dev}); create-react-app has been deprecated by the React team since February 2025.
The tutor is asked to confirm (Appendix~\ref{app:cr}).
\end{decision}

\begin{decision}{D-A3-3: Plotly.js for every chart}
\textbf{Recommendation.} One wrapper component draws every chart with Plotly.js. The library is
loaded only when a chart first appears and is built from the trace types the application uses.

\textbf{Why.} Assignment 1's WBS dictionary named Plotly.js. The brief's High Distinction line
names D3.js or Plotly.js. Zoom, pan, a range slider, hover tooltips, legend toggles and PNG export
are built in, so the effort goes into what the charts say.

\textbf{Alternative.} D3.js gives full control but needs each interaction written by hand;
Chart.js is not on the High Distinction line.
\end{decision}

\begin{decision}{D-A3-4: Axios for every request}
\textbf{Recommendation.} \fp{src/api/http.ts} becomes one Axios instance with a base address, a
time limit, cancellation, and an interceptor that turns every failure into the \code{ApiError} the
screens already handle.

\textbf{Why.} Assignment 1's dictionary says ``call the API with Axios'' and the brief's report
outline gives Axios as its example. Cancellation and time limits are what the live what-if needs.
\end{decision}

\begin{decision}{D-A3-5: Assignment 1's API is kept and extended; validation errors answer 400}
\textbf{Recommendation.} The five planned endpoints keep their paths, methods and bodies. New
endpoints sit under the same \code{/api} prefix. A request that fails validation answers 400 with
the \code{invalid} body.

\textbf{Why.} FastAPI answers 422 for a failed validation by default, but Assignment 1 gave 422
one meaning: the address or transaction has too little history to score. Keeping the two apart lets
the front end map each status to one state.
\end{decision}

\begin{decision}{D-A3-6: the baseline decides the flags; other methods are a view}
\textbf{Recommendation.} The robust-z baseline stays the Explorer's method, as Assignment 2
selected it. A selector (S-01) lets the user see the flags of any of the nine methods; it always
starts on the baseline and says which method is showing.

\textbf{Why.} The models stay frozen (W-06), yet all nine reach the interface, which is what
``AI model integration'' is marked on.
\end{decision}

\begin{decision}{D-A3-7: the Must tier works with no internet}
\textbf{Recommendation.} Everything in the Must tier reads local files. Its real-time behaviour
is the what-if, computed by the server on every change. The replay (S-06) also needs no network.
Only the live lookup (S-09) calls Esplora, behind a setting that is off unless CR-08 is approved.

\textbf{Why.} Assignment 1 risk 15 fires if ``the application needs a live call to run''. A marker
with no connection, or a rate-limited one, still sees every Must feature.
\end{decision}

\begin{decision}{D-A3-8: no accounts; preferences in the browser, the watchlist in SQLite}
\textbf{Recommendation.} Theme, units, recent addresses and the last filters live in the
browser's local storage. The watchlist lives on the server in one SQLite file, written through
Python's standard library. The snapshot stays read-only files.

\textbf{Why.} The brief asks for no login, and the Week 8 lecture says a data layer is not
compulsory. A small store still gives the API its PUT and DELETE and gives the interface something
personal to remember.
\end{decision}

\begin{decision}{D-A3-9: the Assistant is optional, guarded and never streamed}
\textbf{Recommendation.} The template explanation is always shown. The Assistant writes only when
the user asks, its answer is shown only after every guardrail check has passed, and a failed answer
is replaced by the template with the reason. Its weights are an optional download, checked by
SHA-256. Answers are cached, and the answers for the snapshot's flagged transactions are prepared in
advance.

\textbf{Why.} Assignment 2's plan made the Assistant optional and off by default; its report says
the template is authoritative. Showing text token by token would display words the guardrail has not
yet checked.
\end{decision}

\begin{decision}{D-A3-10: one pinned environment per tier}
\textbf{Recommendation.} \code{environment.yml} creates \code{cos30049-a3} with Python 3.11 and
Assignment 2's exact versions of the libraries the saved models were written with: numpy 2.2.6,
pandas 2.2.3, scikit-learn 1.6.1, scipy 1.17.1 and joblib 1.4.2. FastAPI, Uvicorn, Pydantic, httpx
and pytest are pinned to the versions installed on the first day. \code{requirements.txt} is the
pip route. PyTorch 2.7.1 and \code{tokenizers} are only in \code{requirements-assistant.txt}. The
front end needs Node.js 22 or later and installs from \code{package-lock.json}.

\textbf{Why.} A different scikit-learn prints a version warning for every model it loads, and the
rubric's execution line says ``without errors or warnings''. PyTorch is a large download that only
the optional Assistant needs.
\end{decision}

\begin{decision}{D-A3-11: two terminals, one origin}
\textbf{Recommendation.} The README starts Uvicorn on port 8000 and Vite on port 5173. Vite passes
every \code{/api} request to the back end, so the browser sees one origin. The back end also allows
the front end's address by CORS, for anyone who points the client straight at port 8000.

\textbf{Why.} It is the layout the unit teaches (a React server and a FastAPI server), and the
proxy removes the most common first-run failure. Serving the built front end from FastAPI with one
command is Could item C-05.
\end{decision}

\subsection{Change requests}

Assignment 1's change control (its §8.5) needs a change request to add an optional item from the
scope exhibit or to change a must-have item, approved by all four members. Two are raised; their
full text is in Appendix~\ref{app:cr}.

\begin{center}
\small
\begin{tabularx}{\textwidth}{@{}L{1.4cm}Y L{3.6cm}@{}}
\toprule
\textbf{ID} & \textbf{Proposal} & \textbf{If it is rejected} \\
\midrule
CR-07 & Adopt the scope of Section~\ref{sec:scope} as the next scope baseline for WBS 1.4: activate Assignment 1's optional items (a third chart type, zoom and filtering, CSV export), and add the what-if lab, the models page, the method selector, the watchlist, the Assistant and the replay & The Must tier shrinks to Assignment 1's minimum viable demonstration and two charts: the brief's pass level, not its High Distinction list \\
CR-08 & Allow a live lookup of any valid Bitcoin address as a Should item behind a setting. This reverses two Assignment 1 statements: ``no live blockchain calls from the application'' and the exclusion of addresses outside the sample & S-09 and C-06 are dropped. Nothing else depends on them \\
\bottomrule
\end{tabularx}
\end{center}

\subsection{Questions for the tutor}
\label{sec:tutorq}

\begin{enumerate}
  \item Does Assignment 2's arrangement for generative AI (allowed, with a separate declaration)
    apply to Assignment 3, whose brief allows AI for learning, idea generation and debugging and
    asks for all use to be acknowledged in the report?
  \item Is React with TypeScript and Vite acceptable, as in the marked Assignment 1 prototype, in
    place of create-react-app?
  \item Is the due time Sunday 6 December 2026, 23:59 GMT+7, and what are Canvas's size limits
    for the zip and the video?
  \item The models are about 30 MB and go in the zip. The optional 152M Assistant is 610 MB. May
    the README link to it with a SHA-256 instead?
  \item Is running locally from the README the expected demonstration, or should the application
    also be hosted?
  \item May the full API reference sit in the appendix, with at least four endpoints documented
    in the 8 pages?
\end{enumerate}
